% =====================================================================
%  PDX — Photo Diary Exchange · Gesamtspezifikation v1.5
%  Konsolidierte Fassung: v1.3.1 (Basis) + v1.4 (source_album)
%  + v1.4.1 (link) + v1.5 (heading, Markdown-light, quote-source),
%  nomenklatorisch auf dem Stand der Implementierung (hidden -> locked).
%  Basis: PDX_SPEC_v1_4_1.tex (App-Projekt MyPhotoDiary, 02.09.2026);
%  die dortigen "Ausblick"-Punkte sind hier normativ eingearbeitet.
%  Gestaltung streng nach MPD_CI (v1.1): Carlito/Calibri, Titelblatt-
%  Struktur, H1 blau mit Unterlinie, Tabellen D6E4F0/EBF3FB,
%  Code-Bloecke F2F2F2, Inline-Code C_RED. Kein Gold in Dokumenten.
%  Fortgeschrieben 02.09.2026 im MPD-Repo (Server-Chat), docs/specs/.
% =====================================================================
\documentclass[11pt,a4paper]{article}

\usepackage[T1]{fontenc}
\usepackage[utf8]{inputenc}
\usepackage[ngerman]{babel}
\usepackage[sfdefault]{carlito}   % metrisch kompatibel zu Calibri (CI §3)
\usepackage{courier}              % Courier fuer Code (CI §3)
\usepackage{microtype}
\usepackage[table]{xcolor}
\usepackage[a4paper,top=2.0cm,bottom=2.0cm,left=2.5cm,right=2.5cm]{geometry} % CI §4
\usepackage{array}
\usepackage{tabularx}
\usepackage{listings}
\usepackage{titlesec}
\usepackage{fancyhdr}
\usepackage[colorlinks=true]{hyperref}

% -- MPD-CI-Farbpalette (MPD_CI §2) ------------------------------------
\definecolor{mpdnavy}{HTML}{2C3E50}   % C_DARK  — Haupttitel, Zitat
\definecolor{mpdslate}{HTML}{34495E}  % C_MID   — Subtitel
\definecolor{mpdblue}{HTML}{2E75B6}   % C_BLUE  — H1/H2, Links
\definecolor{mpdbluemid}{HTML}{4F81BD}% C_LIGHT — H3/H4, Tabellen-Header
\definecolor{mpdgrey}{HTML}{7F8C8D}   % C_GREY  — Captions, Meta
\definecolor{mpdred}{HTML}{C0392B}    % C_RED   — Inline-Code, Bezeichner
\definecolor{mpdrule}{HTML}{5B9BD5}   % C_LINE  — Trennlinien
\definecolor{tblhead}{HTML}{D6E4F0}   % Tabellen-Kopf
\definecolor{tblalt}{HTML}{EBF3FB}    % alternierende Zeilen
\definecolor{tblborder}{HTML}{CCCCCC} % Tabellenrahmen
\definecolor{codebg}{HTML}{F2F2F2}    % Code-Block-Hintergrund

\hypersetup{linkcolor=mpdblue, urlcolor=mpdblue, citecolor=mpdblue,
  pdftitle={PDX — Photo Diary Exchange, Spezifikation v1.5},
  pdfauthor={Thomas Bugge / My Photo Diary}}

% -- Ueberschriften (CI §3/§4: H1 16pt blau + Unterlinie 5B9BD5) -------
\titleformat{\section}
  {\color{mpdblue}\bfseries\fontsize{16}{20}\selectfont}
  {\thesection.}{0.5em}{}
  [{\color{mpdrule}\titlerule[0.8pt]}]
\titleformat{\subsection}
  {\color{mpdblue}\bfseries\fontsize{13}{16}\selectfont}
  {\thesubsection}{0.5em}{}
\titleformat{\subsubsection}
  {\color{mpdblue}\bfseries\normalsize}{\thesubsubsection}{0.5em}{}

% -- Kopf-/Fusszeile: Meta in C_GREY -----------------------------------
\pagestyle{fancy}\fancyhf{}
\renewcommand{\headrulewidth}{0pt}
\fancyhead[L]{\small\color{mpdgrey}PDX — Photo Diary Exchange}
\fancyhead[R]{\small\color{mpdgrey}Spezifikation v1.5}
\fancyfoot[L]{\small\color{mpdgrey}\itshape where your memories stay yours.}
\fancyfoot[R]{\small\color{mpdgrey}\thepage}

% -- Inline-Code rot (CI §2: C_RED fuer technische Bezeichner) ---------
\let\mpdtt\texttt
\renewcommand{\texttt}[1]{\mpdtt{\textcolor{mpdred}{#1}}}
\newcommand{\feld}[1]{\mpdtt{\textbf{\textcolor{mpdred}{#1}}}}

% -- Code-Bloecke: F2F2F2, kein Rahmen, Courier 9pt (CI §4) ------------
\lstdefinelanguage{json}{
  string=[b]",
  stringstyle=\color{mpdnavy},
  literate=
    {:}{{{\color{mpdgrey}:}}}{1}
    {,}{{{\color{mpdgrey},}}}{1}
    {\{}{{{\color{mpdblue}\{}}}{1}
    {\}}{{{\color{mpdblue}\}}}}{1}
    {[}{{{\color{mpdblue}[}}}{1}
    {]}{{{\color{mpdblue}]}}}{1}
    {ä}{{\"a}}1 {ö}{{\"o}}1 {ü}{{\"u}}1 {ß}{{\ss}}1
    {Ä}{{\"A}}1 {Ö}{{\"O}}1 {Ü}{{\"U}}1,
}
\lstset{
  language=json,
  basicstyle=\ttfamily\footnotesize,
  backgroundcolor=\color{codebg},
  frame=none, xleftmargin=6pt, framexleftmargin=6pt,
  breaklines=true, columns=fullflexible,
  keepspaces=true, showstringspaces=false,
  aboveskip=8pt, belowskip=8pt,
}

% -- Tabellen: Kopf D6E4F0, Zeilen weiss/EBF3FB, Rahmen CCCCCC (CI §4) -
\arrayrulecolor{tblborder}
\setlength{\arrayrulewidth}{0.75pt}
\renewcommand{\arraystretch}{1.3}
\rowcolors{2}{white}{tblalt}
\newcommand{\kopf}[1]{\textbf{#1}}
\newcolumntype{L}{>{\raggedright\arraybackslash}X}
\setlength{\parindent}{0pt}
\setlength{\parskip}{4pt}

\begin{document}
\shorthandoff{"}

% -- Titelblock exakt nach CI §4 (Titelblatt-Struktur) -----------------
\begin{center}
  {\fontsize{36}{42}\selectfont\bfseries\color{mpdnavy} My Photo Diary}\\[6pt]
  {\fontsize{18}{22}\selectfont\color{mpdslate} PDX — Photo Diary Exchange
   \;·\; Format Specification}\\[8pt]
  {\fontsize{13}{16}\selectfont\color{mpdgrey} Version 1.5 \;·\; 2026-09-02}\\[8pt]
  {\color{mpdrule}\rule{\linewidth}{1pt}}\\[6pt]
  {\color{mpdgrey}\small Status: \textbf{Stable} \;·\; konsolidierte
   Gesamtfassung \;·\; Lizenz: Open}\\[6pt]
  {\color{mpdrule}\rule{\linewidth}{1pt}}\\[8pt]
  {\itshape\color{mpdnavy} „My Photo Diary — where your memories stay yours.“}
\end{center}

\vspace{4pt}
{\small\color{mpdgrey}
Diese Gesamtfassung konsolidiert PDX v1.3.1 (Elementtypen), v1.4
(\texttt{source\_album}), v1.4.1 (\texttt{link}) und v1.5
(\texttt{heading}, Markdown-light, \texttt{source} am Zitat) in einem
Dokument und
verwendet durchgehend die aktuelle Nomenklatur der MPD-Implementierung
(\texttt{locked}/\texttt{show\_locked}; siehe Abschnitt~\ref{sec:locked}).
Referenzimplementierung: MPD~v2 — \texttt{backend/routers/albums.py}
(Validierung, Auflösung), \texttt{frontend/js/album-render.js}
(Darstellung) sowie die native Android-App (\texttt{de.bgghome.mpd}).}

\tableofcontents

% ======================================================================
\section{Konzept}

PDX beschreibt ein Fotoalbum als \emph{Drehbuch}: eine einzige, von Hand
lesbare Datei \texttt{album.json} neben den unveränderten Mediendateien.
Die Grundsätze:

\begin{itemize}
  \item \textbf{Overlay-Prinzip.} PDX verändert nie Mediendateien. Das
        Löschen eines Albums löscht nur das Drehbuch; die Photos bleiben.
  \item \textbf{Verweis-Container.} Jedes Medien-Element ist ein Verweis
        auf eine Datei — keine eingebetteten Daten. Seit v1.4 kann der
        Verweis auch in einen anderen Albumordner desselben Spaces zeigen.
  \item \textbf{Reihenfolge ist Wahrheit.} Die Array-Reihenfolge in
        \texttt{elements} ist die einzige Quelle der Anzeige-Reihenfolge.
  \item \textbf{Vorwärtskompatibilität.} Viewer \emph{müssen} unbekannte
        Elementtypen und Felder still überspringen — nie abstürzen.
  \item \textbf{Deklarativ, nicht gerendert.} PDX speichert Stil-\emph{Namen}
        (z.\,B. \texttt{"style": "quote"}); wie diese aussehen, entscheidet
        jeder Viewer (Web-CSS, native App, künftige Ausgaben wie ein
        PDF-Fotobuch) — mit dem Ziel eines konsistenten Gesamtbilds.
        Die \emph{Referenzdarstellung} der Stile (Schrift, Größe, Farbe,
        Form) legt die MPD-CI fest (\texttt{MPD\_CI\_v1\_2}, Abschnitt
        „Referenz-Typografie der PDX-Textstile“); eigene Viewer sollten
        sich daran orientieren, müssen es aber nicht.
\end{itemize}

\subsection{Dateistruktur}

Flaches Verzeichnis, keine Unterordner (abgesehen vom verwalteten
Sicherungsordner der Implementierung):

\begin{lstlisting}[language={}]
2024-gran-canaria/
  album.json
  2024-03-15_08-30-15.jpg        <- photo
  2024-03-15_16-45-12.mp4        <- video
  2024-03-15_hike.gpx            <- tour
  2024-03-16_hotelrechnung.pdf   <- document (PDF)
  2024-03-15_gipfel.m4a          <- audio
  .mpd_backups/                  <- MPD: letzte 10 Staende von album.json
\end{lstlisting}

% ======================================================================
\section{Wurzelstruktur und Meta}

\begin{lstlisting}
{
  "version": "1.5",
  "meta": { ... },
  "elements": [ ... ]
}
\end{lstlisting}

\subsection{\texttt{meta}}

\begin{tabularx}{\linewidth}{|l|l|l|L|}
\hline
\rowcolor{tblhead}\kopf{Feld} & \kopf{Typ} & \kopf{Pflicht} & \kopf{Bedeutung} \\
\hline
\feld{title} & String & ja & Anzeigetitel des Albums \\
\feld{description} & String & nein & Kurzbeschreibung (Übersichtskarte) \\
\feld{year} & Zahl & nein & Jahr; \texttt{null} bei jahrübergreifenden Sammlungen \\
\feld{thumbnail} & String & nein & Dateiname des Covers, immer lokal im Albumordner \\
\feld{created} & String & ja & ISO-8601-Zeitstempel der Erstellung \\
\feld{modified} & String & ja & ISO-8601; wird vom Server bei jedem Speichern gesetzt \\
\feld{show\_locked} & Boolean & nein & Gesperrte Elemente im Lesemodus zeigen? Default \texttt{true} (Abschnitt~\ref{sec:locked}) \\
\hline
\end{tabularx}

\emph{MPD-Laufzeitfelder} (nicht Teil der Spezifikation, von der
Implementierung ergänzt): \texttt{thumbnail\_crop} (Cover-Ausschnitt der
Übersicht).

% ======================================================================
\section{Elementtypen}

Gemeinsame Basisfelder aller Elemente:

\begin{tabularx}{\linewidth}{|l|l|l|L|}
\hline
\rowcolor{tblhead}\kopf{Feld} & \kopf{Typ} & \kopf{Pflicht} & \kopf{Bedeutung} \\
\hline
\feld{id} & String & ja & vierstellig (\texttt{0001}–\texttt{9999}), albumweit eindeutig, stabil \\
\feld{type} & String & ja & einer der Typen dieses Abschnitts \\
\feld{locked} & Boolean & nein & Element privat: aus geteilten Alben ausgeblendet; im Lesemodus nur sichtbar, wenn \texttt{meta.show\_locked} es erlaubt \\
\hline
\end{tabularx}

\subsection{\texttt{photo}}

\begin{lstlisting}
{ "id": "0005", "type": "photo", "file": "2024-03-15_08-30-15.jpg" }
\end{lstlisting}

\begin{tabularx}{\linewidth}{|l|l|l|L|}
\hline
\rowcolor{tblhead}\kopf{Feld} & \kopf{Typ} & \kopf{Pflicht} & \kopf{Bedeutung} \\
\hline
\feld{file} & String & ja & Dateiname (JPG/PNG/WebP …), relativ zum Albumordner \\
\feld{source\_album} & String & nein & Auflösung aus anderem Albumordner (Abschnitt~\ref{sec:source}) \\
\hline
\end{tabularx}

Bildunterschriften sind \emph{kein} PDX-Feld: Der Erzählfluss lebt von
\texttt{text}-Elementen \emph{zwischen} den Photos (bewusste
Format-Entscheidung); eine dateigebundene Beschriftung gehört ins EXIF
der Bilddatei (\texttt{ImageDescription}).

\subsection{\texttt{video}}

Felder wie \texttt{photo}. Darstellung: statisches Vorschaubild im
Lesefluss; Antippen öffnet einen Player (HTML5-\texttt{<video>} bzw.
nativer Player), ohne Autoplay.

\subsection{\texttt{text}}

\begin{lstlisting}
{ "id": "0002", "type": "text",
  "text": "Endlich auf der Insel.", "style": "default" }
\end{lstlisting}

\begin{tabularx}{\linewidth}{|l|l|l|L|}
\hline
\rowcolor{tblhead}\kopf{Feld} & \kopf{Typ} & \kopf{Pflicht} & \kopf{Bedeutung} \\
\hline
\feld{text} & String & ja & Text; Zeilenumbrüche erlaubt, kein HTML — seit v1.5 mit Markdown-light (Abschnitt~\ref{sec:mdlight}) \\
\feld{style} & String & nein & \texttt{default} · \texttt{info} · \texttt{quote} · \texttt{note} · \texttt{heading} \emph{(v1.5)} \\
\feld{source} & String & nein & \emph{v1.5:} Quellenangabe, nur bei \texttt{style: "quote"} definiert; reiner Text, Anzeige klein unter dem Zitat \\
\hline
\end{tabularx}

\begin{tabularx}{\linewidth}{|l|L|}
\hline
\rowcolor{tblhead}\kopf{Stil} & \kopf{Charakter} \\
\hline
\feld{default} & Erzähltext, zentriert, lesefreundliche Breite \\
\feld{info} & sachlicher Kasten, blaue Akzentkante \\
\feld{quote} & Zitat, kursiv, mit Anführungs-Initiale \\
\feld{note} & persönliche Notiz, goldene Akzentkante \\
\feld{heading} & \emph{v1.5:} echte Kapitel-Überschrift — größere Schrift, Akzentfarbe des Themes; ersetzt den Behelf über Trenner-Labels \\
\hline
\end{tabularx}

\subsubsection*{Markdown-light \;(v1.5)}\label{sec:mdlight}

Im \feld{text}-Feld sind genau \emph{zwei} Auszeichnungen zulässig:
\texttt{*kursiv*} und \texttt{**fett**}. Kein HTML, keine weiteren
Markdown-Konstrukte, keine Verschachtelung.

\begin{itemize}
  \item \textbf{Speicherung:} \feld{text} bleibt ein reiner String.
        Server \textbf{dürfen} nichts strippen, rendern oder
        umschreiben — die Sternchen sind Teil des gespeicherten Texts.
  \item \textbf{Darstellung:} Viewer rendern die beiden Auszeichnungen
        in allen Text-Stilen. Die Umsetzung \textbf{muss} injektionssicher
        sein (Referenz \texttt{album-render.js}: Tokenisierung in Text-/%
        \texttt{<strong>}/\texttt{<em>}-Knoten, kein \texttt{innerHTML}).
  \item \textbf{Editoren} geben beim Bearbeiten den Roh-Text samt
        Sternchen aus — gerendertes Markup im Editor würde die
        Auszeichnung beim Zurückschreiben zerstören.
  \item \textbf{Degradation:} Viewer vor v1.5 zeigen die Sternchen
        wörtlich — die einzige, bewusst verschmerzte Abweichung.
\end{itemize}

\subsection{\texttt{separator}}

Erzählerische Zäsur — Kapitelwechsel, Zeitsprung.

\begin{lstlisting}
{ "id": "0015", "type": "separator", "label": "Tag 2 - Florenz" }
\end{lstlisting}

\begin{tabularx}{\linewidth}{|l|l|l|L|}
\hline
\rowcolor{tblhead}\kopf{Feld} & \kopf{Typ} & \kopf{Pflicht} & \kopf{Bedeutung} \\
\hline
\feld{label} & String & nein & kurze Beschriftung, max.\ 80 Zeichen; ohne Label reine Trennlinie \\
\feld{style} & String & nein & \texttt{line} (Default) · \texttt{bold} · \texttt{dashed} · \texttt{dots} · \texttt{space} \\
\hline
\end{tabularx}

\subsection{\texttt{tour}}

Lokal gespeicherter GPS-Track (GPX), als interaktive Karte gerendert —
vollständig selbst gehostet.

\begin{lstlisting}
{ "id": "0016", "type": "tour", "file": "2024-03-15_hike.gpx",
  "title": "Roque Nublo", "source": "komoot" }
\end{lstlisting}

\begin{tabularx}{\linewidth}{|l|l|l|L|}
\hline
\rowcolor{tblhead}\kopf{Feld} & \kopf{Typ} & \kopf{Pflicht} & \kopf{Bedeutung} \\
\hline
\feld{file} & String & ja & GPX-Datei im Albumordner \\
\feld{title} & String & nein & Anzeigetitel über der Karte \\
\feld{source} & String & nein & Herkunft (\texttt{komoot}, \texttt{garmin} …), rein informativ \\
\feld{type\_hint} & String & nein & Sportart für das Icon (\texttt{hiking}, \texttt{cycling} …) \\
\feld{map\_layer} & String & nein & Grundkarte \texttt{osm} (Default) · \texttt{topo} \\
\feld{show\_elevation} & Boolean & nein & Höhenprofil anzeigen, Default \texttt{true} \\
\feld{source\_album} & String & nein & siehe Abschnitt~\ref{sec:source} \\
\hline
\end{tabularx}

\subsection{\texttt{document}}

Gescannte Belege — Tickets, Rechnungen, Bordkarten — und seit 1.5.1 auch
Textdateien. \emph{PDF} wird als Karte mit Reader dargestellt;
\emph{JPG/PNG-Scans} werden inline wie Photos gesetzt; \emph{TXT/MD}
erscheinen als aufklappbarer Kasten, dessen Inhalt beim Aufklappen aus der
Datei gelesen wird — nicht in die \texttt{album.json} kopiert, damit es nur
eine Wahrheit gibt. Der Text ist auf 100\,KB gedeckelt, darüber verweist
der Kasten auf die vollständige Ansicht.

\feld{title} ist bei \texttt{document} zugleich die Beschriftung
\emph{im} Element: in der PDF-Karte an Stelle des Dateinamens, bei einer
Textdatei in der Klappzeile. Fehlt er, zeigt die Darstellung eine
abgeleitete Beschriftung (bei Textdateien die erste nicht-leere Zeile,
sonst den Dateinamen). Diese abgeleitete Beschriftung wird \emph{nicht}
gespeichert — in der \texttt{album.json} steht dann weiterhin kein
\feld{title}.

\begin{lstlisting}
{ "id": "0021", "type": "document",
  "file": "2024-03-16_hotelrechnung.pdf",
  "title": "Hotel Parador - Rechnung", "kind": "receipt" }
\end{lstlisting}

\begin{tabularx}{\linewidth}{|l|l|l|L|}
\hline
\rowcolor{tblhead}\kopf{Feld} & \kopf{Typ} & \kopf{Pflicht} & \kopf{Bedeutung} \\
\hline
\feld{file} & String & ja & JPG/PNG/PDF/TXT/MD im Albumordner \\
\feld{title} & String & nein & Beschriftung im Element (PDF-Karte, Klappzeile); bei JPG/PNG darüber \\
\feld{kind} & String & nein & \texttt{ticket} · \texttt{receipt} · \texttt{invoice} · \texttt{boarding\_pass} · \texttt{other}; rein informativ \\
\feld{source\_album} & String & nein & siehe Abschnitt~\ref{sec:source} \\
\hline
\end{tabularx}

\subsection{\texttt{audio}}

Sprachnotizen und Atmosphären, als Inline-Player gerendert.

\begin{lstlisting}
{ "id": "0022", "type": "audio",
  "file": "2024-03-15_gipfel.m4a", "title": "Wind am Gipfel" }
\end{lstlisting}

Felder: \feld{file} (Pflicht; MP3/M4A/OGG/WAV), \feld{title} (optional),
\feld{source\_album} (optional).

\subsection{\texttt{map}}

Manuell gesetzte Orts-Pins ohne Track — Reiserouten, Hotels, Sehenswertes.

\begin{lstlisting}
{ "id": "0023", "type": "map", "title": "Reiseroute",
  "pins": [
    { "lat": 53.6304, "lon":  9.9882, "label": "Hamburg (HAM)" },
    { "lat": 27.9319, "lon": -15.3866, "label": "Gran Canaria (LPA)" }
  ] }
\end{lstlisting}

\begin{tabularx}{\linewidth}{|l|l|l|L|}
\hline
\rowcolor{tblhead}\kopf{Feld} & \kopf{Typ} & \kopf{Pflicht} & \kopf{Bedeutung} \\
\hline
\feld{pins} & Array & ja & mindestens ein Pin: \feld{lat}/\feld{lon}
  (WGS84, Pflicht), \feld{label} (optional) \\
\feld{title} & String & nein & Anzeigetitel über der Karte \\
\feld{layer} & String & nein & Grundkarte \texttt{osm} (Default) · \texttt{topo} \\
\hline
\end{tabularx}

Darstellung: Zoom automatisch auf alle Pins; ein Pin → nur Marker;
ab zwei Pins verbindet eine Geodäten-Linie (Großkreis) die Pins in
Reihenfolge — ideal für Reiserouten.

\subsection{\texttt{link} \; (neu in v1.4.1)}\label{sec:link}

Verweis auf eine Webseite — der erste Elementtyp, der \emph{nicht} auf
eine Datei zeigt: statt \feld{file} trägt er \feld{url}.

\begin{lstlisting}
{ "id": "0042", "type": "link",
  "url": "https://www.alpenvereinaktiv.com/de/tour/12345",
  "title": "Wanderung Rifugio Gardeccia",
  "note": "Route und Hoehenprofil" }
\end{lstlisting}

\begin{tabularx}{\linewidth}{|l|l|l|L|}
\hline
\rowcolor{tblhead}\kopf{Feld} & \kopf{Typ} & \kopf{Pflicht} & \kopf{Bedeutung} \\
\hline
\feld{url} & String & \textbf{ja} & nur \texttt{http://}/\texttt{https://}, max.\ 2048 Zeichen \\
\feld{title} & String & nein & Anzeigename; fehlt er, zeigt der Viewer die Domain \\
\feld{note} & String & nein & eine Zeile Kontext unter dem Chip \\
\hline
\end{tabularx}

Regeln (Sicherheit und Privatsphäre):
\begin{itemize}
  \item Andere Schemata (\texttt{javascript:}, \texttt{data:},
        \texttt{file:} …) sind \textbf{unzulässig}; die Prüfung
        \textbf{muss} serverseitig erfolgen.
  \item Der Viewer \textbf{muss} in neuem Kontext öffnen und \textbf{soll}
        die Zieldomain vor dem Aktivieren sichtbar machen.
  \item Weder Client noch Server \textbf{dürfen} die Zielseite für
        Vorschaudaten abrufen — in geteilten Alben würde das dem
        Zielserver Existenz und Zeitpunkt des Betrachtens verraten.
        Titel und Notiz werden von Hand gepflegt.
  \item \feld{source\_album} ist für \texttt{link} nicht definiert.
\end{itemize}

% ======================================================================
\section{Verweise in andere Alben: \texttt{source\_album} \;(v1.4)}\label{sec:source}

Medien-Elemente (\texttt{photo}, \texttt{video}, \texttt{tour},
\texttt{document}, \texttt{audio}) dürfen optional
\feld{source\_album} tragen: Die Datei wird dann im genannten
Albumordner \emph{desselben Spaces} aufgelöst statt lokal.

\begin{lstlisting}
{ "id": "0042", "type": "photo", "file": "IMG_5678.jpg",
  "source_album": "2022-einschulung" }
\end{lstlisting}

\begin{itemize}
  \item Anwendungsfall: kuratierte Sammlungen (Personen-Alben über Jahre,
        Best-ofs) ohne Datei-Duplikate. Mischformen sind erlaubt; es gibt
        keinen eigenen Album-Typ.
  \item \textbf{Cross-Space-Referenzen sind verboten} (Space-Isolation,
        schlichtes Rechte-Modell). Validatoren müssen sie ablehnen.
  \item Verlorene Referenzen führen zu Soft-Fail
        (\texttt{source\_missing: true} zur Laufzeit), nie zum Absturz.
  \item \textbf{Cover:} \texttt{meta.thumbnail} bleibt immer ein lokaler
        Dateiname. Ein referenziertes Photo wird als Cover physisch als
        \texttt{cover.<ext>} in den eigenen Ordner kopiert
        (MPD: \texttt{POST \dots/cover-from-source}).
\end{itemize}

% ======================================================================
\section{Sichtbarkeit: \texttt{locked} und \texttt{show\_locked}}\label{sec:locked}

Jedes Element kann mit \texttt{locked: true} als privat markiert werden:

\begin{itemize}
  \item In \textbf{geteilten} Alben (Share-Links) werden gesperrte
        Elemente immer ausgefiltert.
  \item Im \textbf{Lesemodus} des Eigentümers entscheidet
        \texttt{meta.show\_locked} (Default \texttt{true}); der
        Schlüssel-Umschalter der Oberfläche dient als Vorschau
        „so sieht es ein Empfänger“.
  \item Im \textbf{Bearbeitungsmodus} sind gesperrte Elemente immer
        sichtbar — sonst wären sie nach dem Sperren unerreichbar.
  \item \texttt{separator}-Elemente sind Struktur, keine Inhalte —
        sie unterliegen dieser Filterung nicht.
\end{itemize}

\emph{Historie:} Bis v1.3.1 hießen diese Konzepte \texttt{hidden} /
\texttt{show\_hidden} (der Text-Stil \texttt{hidden} eingeschlossen).
Die Implementierung migriert Altbestände beim Laden still auf
\texttt{locked} / \texttt{show\_locked}.

% ======================================================================
\section{Laufzeitfelder (nicht speichern)}

Die Auslieferung ergänzt berechnete Felder, die \textbf{nie} in
\texttt{album.json} zurückgeschrieben werden dürfen; Server
\textbf{müssen} sie beim Speichern verwerfen:

\begin{tabularx}{\linewidth}{|l|l|L|}
\hline
\rowcolor{tblhead}\kopf{Feld} & \kopf{an Typen} & \kopf{Bedeutung} \\
\hline
\feld{thumbnails} & photo, video, document & URLs der Größen \texttt{sm}/\texttt{m}/\texttt{xl}/\texttt{full} \\
\feld{original} & photo & Download-URL des Originals \\
\feld{resolution} & photo & \texttt{\{width, height\}} in Pixeln \\
\feld{blurhash} & photo & Platzhalter-Hash für sanftes Laden \\
\feld{document\_url} & document (PDF) & Auslieferungs-URL des Readers \\
\feld{error} & Medien-Typen & gesetzt, wenn die Datei fehlt (Soft-Fail) \\
\feld{source\_missing} & Medien-Typen & \texttt{true}, wenn eine \texttt{source\_album}-Referenz ins Leere zeigt \\
\hline
\end{tabularx}

% ======================================================================
\section{Validierung (Zusammenfassung)}

\begin{itemize}
  \item \feld{id}: Pflicht, vierstellig, albumweit eindeutig.
  \item Alle Dateiverweise müssen auf existierende Dateien zeigen
        (Soft-Fail zur Laufzeit, harte Prüfung beim Anlegen).
  \item \texttt{separator.label} max.\ 80 Zeichen, reiner Text.
  \item \texttt{map.pins} mindestens ein Pin mit \feld{lat}/\feld{lon}.
  \item \texttt{document.file}: JPG, PNG, PDF, TXT oder MD; \texttt{audio.file}:
        MP3, M4A, OGG oder WAV; \texttt{tour.file}: gültiges GPX.
  \item \texttt{link.url}: Pflicht, nur \texttt{http}/\texttt{https},
        max.\ 2048 Zeichen; Prüfung serverseitig.
  \item \texttt{source\_album}: einfacher Ordnername ohne Pfadtrenner
        und ohne \texttt{..}; gleicher Space; Cross-Space ablehnen.
  \item Titel/Notizen/Labels sind reiner Text — Viewer setzen sie als
        Text, nie als Markup. Gleiches gilt für \texttt{text.source}
        (v1.5). Markdown-light gilt \emph{ausschließlich} für
        \feld{text} (Abschnitt~\ref{sec:mdlight}).
\end{itemize}

% ======================================================================
\section{Kompatibilität und Migration}

Jede Version seit v1.2.1 ist ein striktes Superset ihrer Vorgängerin —
\textbf{keine Migration nötig}, alte Alben bleiben gültig. Viewer
überspringen Unbekanntes still. \texttt{version} \textbf{soll} auf die
höchste genutzte Fähigkeit gesetzt werden (\texttt{"1.4.1"}, sobald
\texttt{link} vorkommt; \texttt{"1.5"}, sobald \texttt{heading} oder
\texttt{source} vorkommt). Die MPD-Referenzimplementierung leitet die
Nummer beim Speichern aus dem Inhalt ab; reines Markdown-light ist im
String nicht sicher erkennbar und ändert die Nummer allein nicht —
unkritisch, da ältere Viewer die Sternchen nur wörtlich zeigen.

\subsection*{Versions-Historie}

\begin{tabularx}{\linewidth}{|l|l|L|}
\hline
\rowcolor{tblhead}\kopf{Version} & \kopf{Datum} & \kopf{Kern} \\
\hline
1.0 & 2025-01-27 & Erststand: \texttt{text} und \texttt{photo} \\
1.1 & 2025-02-25 & Feldnamen englisch; ID-System; Text-Stile \\
1.2 & 2025-03-14 & \feld{id} Pflicht; Array-Reihenfolge als einzige Wahrheit \\
1.2.1 & 2025-03-14 & Text-Stile wiederhergestellt; Vollbeispiel \\
1.3 & 2025-04-04 & \texttt{separator}, \texttt{tour}, \texttt{document}, \texttt{audio}, \texttt{map}; Vorwärtskompatibilität verbindlich \\
1.3.1 & 2025-04-06 & Klarstellungen (PDF-Dokumente, Geodäten-Linie) \\
1.4 & 2026-04-21 & \feld{source\_album}, \feld{source\_missing}; Cover-Kopie-Regel \\
1.4.1 & 2026-08-23 & \texttt{link}-Element (Abschnitt~\ref{sec:link}) \\
1.5 & 2026-09-02 & Text-Stil \texttt{heading}; Markdown-light
  (Abschnitt~\ref{sec:mdlight}); \feld{source} am Zitat \\
1.5.1 & 2026-09-15 & \texttt{document}: TXT/MD als aufklappbarer Kasten;
  \feld{title} steht im Element statt darüber \\
\hline
\end{tabularx}

\emph{Anmerkung zur Nummer 1.4.1:} Nach der Versionsregel der MPD-CI
(Patch-Stellen nur für Errata, neue Fähigkeiten zählen MINOR) hätte das
\texttt{link}-Element eigentlich \textbf{v1.5} heißen müssen. Die Nummer
1.4.1 ist jedoch bereits in Server, Repo und gespeicherten Alben
eingeführt; dieses Dokument behält sie bei. Diese Fassung zählt wieder
regelkonform: Die Textgestaltungs-Erweiterung ist ein Feature und damit
MINOR — \textbf{v1.5}.

% ======================================================================
\section{Ausblick}\label{sec:ausblick}

Die im v1.4.1-Dokument als Entwurf geführten Punkte (\texttt{heading},
Markdown-light, \texttt{source}) sind in dieser Fassung normativ.
Weiterhin gilt: \emph{keine} Bildunterschriften in PDX — dateigebundene
Beschriftungen gehören ins EXIF (\texttt{ImageDescription}); ein
künftiges Fotobuch-PDF kann sie von dort als klassische
Bildunterschriften setzen.

% ======================================================================
\section{Lizenz}

Diese Spezifikation ist offen und frei nutzbar — keine Einschränkungen
für Implementierung oder Verwendung. Entworfen für die langfristige
Archivierung persönlicher Foto-Sammlungen in einem menschenlesbaren,
zukunftssicheren Format.

\end{document}
